% Generated from es/cv.md by tools/build_cv.py. Do not edit.



\section{Puestos académicos}

\begin{entries}
\entry{2026–}{Profesor de planta de Economía, Escuela de Ciencias Sociales y Gobierno, Tecnológico de Monterrey}
\entry{2025–2026}{Investigador posdoctoral y profesor de cátedra, Tecnológico de Monterrey}
\entry{2024–2025}{Investigador posdoctoral y profesor asistente temporal (Term Assistant Professor), George Mason University Korea}
\end{entries}

\section{Formación}

\begin{entries}
\entry{2023}{Doctorado en Economía, George Mason University (Interdisciplinary Center for Economic Science)}
\entry{2013}{Maestría en Economía, Universidad de los Andes}
\entry{2012}{Pregrado en Economía, Universidad de los Andes}
\end{entries}

\section{Áreas}

Microeconomía aplicada, economía experimental, economía pública, economía política.

\section{Publicaciones}

\begin{entries}
\entry{}{``Public Good Provision and Property Tax Compliance: Evidence from a Natural Experiment'' (con Paul E. Carrillo y Carlos Scartascini). \emph{Journal of Public Economics}, 198 (2021): 104422. \href{https://doi.org/10.1016/j.jpubeco.2021.104422}{DOI}}
\end{entries}

\section{Documentos de trabajo}

\begin{entries}
\entry{}{``Tackling Biased Beliefs to Foster Women's Entry into Economics: Experimental Evidence from Mexico'' (con Stanislao Maldonado, Elvira Naranjo y Joana Chapa). Documento de trabajo, 2026. \href{https://www.stanislaomaldonado.org/documents/research/tackling-biased-beliefs.pdf}{PDF}}
\entry{}{``Price Distortions and Hoarding: An Experiment'' (con César Martinelli). GMU Working Paper in Economics No. 23-45, 2023. Revise and resubmit, \emph{Journal of Economic Behavior \& Organization}. \href{https://ssrn.com/abstract=4654570}{SSRN}}
\entry{}{``Discretionary Procurement Flexibilization, Efficiency, and Rent-Seeking: Evidence from Chile During COVID-19'' (con Anastasiya Yarygina-Udovenko y Gabriel Villalba Ortega). Serie de Documentos de Trabajo del BID, 2026. \href{https://publications.iadb.org/en/discretionary-procurement-flexibilization-efficiency-and-rent-seeking-evidence-chile-during-covid}{Enlace}}
\entry{}{``Imperfect Attention in Public Policy: A Field Experiment during a Tax Amnesty in Argentina'' (con Carlos Scartascini). Documento para discusión del BID No. IDB-DP-665, 2019. \href{https://doi.org/10.18235/0001661}{DOI}}
\end{entries}

\section{Trabajo en curso}

\begin{entries}
\entry{}{``Beliefs About Price Controls and Zero-Sum Thinking'' (con Stanislao Maldonado y Santiago Sautua)}
\entry{}{``Fairness, Misinformation, and Tax Policy Preferences'' (con Stanislao Maldonado)}
\entry{}{``Affirmative Action, Sorting, and Efficiency in Competition for Selective Colleges'' (con Stanislao Maldonado y César Martinelli)}
\entry{}{``Discretionality of Public Procurement and Efficiency: Colombia'' (con Anastasiya Yarygina-Udovenko)}
\entry{}{``Patronage Networks and Favoritism in Colombian Public Procurement'' (con Hugo Díaz)}
\entry{}{``Collective Deliberation: Theory and Experiment'' (con Mikhail Freer y César Martinelli)}
\end{entries}

\section{Presentaciones en congresos}

\begin{entries}
\entry{2026}{Academics Against Corruption in the Americas, Quito; Public Choice Society Annual Meeting, San Antonio}
\entry{2025}{LACEA Annual Meeting, Recife; Workshop on Behavioral and Experimental Economics, SUFE, Shanghái; APEE Annual Meeting, Antigua Guatemala}
\entry{2022}{APEE Annual Meeting, Las Vegas}
\entry{2021}{APEE Annual Meeting, Fort Lauderdale; Shadow Economy, Tax Behavior, and Institutions Workshop (en línea)}
\entry{2019}{National Tax Association Annual Meeting, Tampa}
\end{entries}

\section{Becas y distinciones}

\begin{entries}
\entry{2024, 2025}{Faculty Research and Development Award, George Mason University Korea (otoño 2024, primavera 2025)}
\entry{2024}{Beca de viaje de la Society for the Advancement of Economic Theory (SAET), Santiago (otorgada; no asistí)}
\entry{2021, 2022}{Young Scholar Program Award, Association of Private Enterprise Education}
\entry{2017–2023}{Beca de asistente de investigación de posgrado, ICES, George Mason University}
\entry{2007–2011}{Beca completa de pregrado ``Quiero Estudiar'', Universidad de los Andes}
\entry{2006}{Mejor puntaje en el examen nacional de admisión a la universidad (ICFES), puesto 1/1000 a nivel nacional}
\end{entries}

\section{Docencia}

\group{\textbf{Tecnológico de Monterrey}}

\begin{entries}
\entry{2026}{EC1025 De microincentivos a macroresultados}
\entry{2026}{EC1015 Economía conductual}
\entry{2026}{EO 4024 Precios y Comportamiento de los Agentes}
\end{entries}

\group{\textbf{George Mason University Korea}}

\begin{entries}
\entry{2024–2025}{ECON 103 Principles of Microeconomics}
\entry{2024–2025}{ECON 306 Intermediate Microeconomics}
\entry{2025}{ECON 495 Policy Impact Evaluation}
\entry{2024}{ECON 104 Principles of Macroeconomics; ECON 311 Intermediate Macroeconomics; ECON 445 Design and Analysis of Experiments}
\end{entries}

\group{\textbf{George Mason University}}

\begin{entries}
\entry{2022}{Instructor, Graduate Math Camp (curso propedéutico de matemáticas para posgrado)}
\entry{2021–2022}{Asistente de docencia, ECON 886 Experimental Economics II (teoría de juegos, doctorado)}
\entry{2021}{Asistente de docencia, ECON 535 Applied Econometrics}
\entry{2020–2021}{Asistente de docencia, ECON 811 Microeconomic Theory II (doctorado)}
\end{entries}

\group{\textbf{Universidad de los Andes}}

\begin{entries}
\entry{2013–2014}{Profesor, Introducción a la Microeconomía}
\entry{2012–2013}{Asistente de docencia, Regulación y Política de Competencia (posgrado)}
\end{entries}

\section{Experiencia en investigación y profesional}

\begin{entries}
\entry{2017–2023}{Asistente de investigación de posgrado, George Mason University}
\entry{2022}{Consultor, Banco Interamericano de Desarrollo (IFD/FMM-0007)}
\entry{2021}{Consultor, Banco Interamericano de Desarrollo (RES-0006)}
\entry{2018–2019}{Investigador predoctoral, CESED, Universidad de los Andes}
\entry{2018}{Consultor, Banco Interamericano de Desarrollo (RES-0005)}
\entry{2014–2017}{Investigador (Research Fellow), Banco Interamericano de Desarrollo, Washington, D.C.}
\entry{2014}{Consultor externo, SENA, Bogotá}
\entry{2011–2014}{Investigador (practicante en 2011), ECON ESTUDIO, Bogotá}
\end{entries}

\section{Actividades profesionales}

\textbf{Dictaminador:} European Economic Review; Games and Economic Behavior; Journal of Economic Behavior \& Organization; Journal of Public Economics.

\textbf{Asesoría:} Codirector de tesis, Universidad del Pacífico, Lima (Matías Villalba), 2023.

\textbf{Escuelas de verano:} ISEO Summer School of Economics, Iseo, Italia, 2023.

\section{Otros}

\textbf{Software:} Stata, R, Python, oTree.\\
\textbf{Idiomas:} español (nativo), inglés (fluido).
